% Q4a in the program's own sense: select the holographic pair.
% Selection-uniqueness is answered in the negative.
% Existence of some pair remains [O]. No CFT is invented.
\documentclass[11pt]{article}
\usepackage[margin=1.05in]{geometry}
\usepackage{amsmath,amssymb,amsthm}
\usepackage{enumitem,booktabs}
\usepackage[colorlinks=true,linkcolor=blue,citecolor=blue,urlcolor=blue]{hyperref}

\newtheorem{theorem}{Theorem}[section]
\newtheorem{proposition}[theorem]{Proposition}
\theoremstyle{definition}
\newtheorem{definition}[theorem]{Definition}
\newtheorem*{admission}{Admission}

\newcommand{\pP}{\textbf{[P]}}
\newcommand{\pC}{\textbf{[C]}}
\newcommand{\pO}{\textbf{[O]}}

\title{\textbf{Q4a in the Proper Context:\\
The Holographic Pair is Not Selected}}
\author{Robert W.\ McGwier\\
\small CTO, Cohere Technology Group\\
\small GitHub: \texttt{n4hy}}
\date{15 August 2026 --- Version 1}

\begin{document}
\maketitle

\begin{abstract}
\noindent
Paper~11 left ultraviolet completion open by construction and refused
to invent a CFT. That refusal was correct as a ban on fabrication.
It is not the end of the question the program actually asked.

The monograph states the question precisely: the New Standard Model
is emergent bookkeeping; the microscopic description is the boundary
theory of a holographic pair; no first-principles selection of that
pair exists \cite{McGwierMonograph}. That is Q4a. It is a
\emph{selection} problem, not an invitation to name $\mathcal{N}=4$
super-Yang--Mills or a Calabi--Yau.

This note answers the selection problem. The certified principles
(Casini--Huerta--Myers, Faulkner--Guica--Hartman--Myers--Van
Raamsdonk, the Paper~II modular rule) determine the bulk
\emph{gravitational} sector given a pair. They are blind to
$G_{\mathrm{SM}}$, the number of generations, and the Yukawas.
They therefore underdetermine the pair. Independently, the Standard
Model cannot \emph{be} the boundary CFT: its one-loop coefficients
are
$(b_1,b_2,b_3)=(41/10,-19/6,-7)$, all nonzero, and it has $118$
light on-shell degrees of freedom, not a large-$N$ gap
\cite{Heemskerk2009}.

Existence of some other pair whose bulk infrared is the NSM remains
\pO{}. No pair is constructed. Final Status ``UV completion'' is
not moved to \pP{}.
\end{abstract}

\section*{Standard of proof}

Q4a, in the vocabulary of this program, is the following pair of
questions.

\begin{enumerate}[leftmargin=1.6em]
\item \textbf{Selection.} Do the certified entanglement principles
single out a holographic pair $(\mathrm{CFT},\mathrm{bulk})$ of
which \eqref{eq:nsm} is the bookkeeping?
\item \textbf{Existence.} Does at least one such pair exist?
\end{enumerate}

A named CFT offered without a derivation from those principles is a
fabrication \cite{McGwierXI}. A uniqueness theorem, or a proof that
uniqueness fails, is an answer to (1). Item (2) is not solved by
solving (1).

\section{What the pair is}

The NSM action is the bulk effective theory
\begin{equation}
S_{\mathrm{NSM}}
=
\frac{1}{4\kappa}
\int\varepsilon_{abcd}\,e^a\wedge e^b\wedge R^{cd}
-
\frac{\Lambda}{6}
\int\varepsilon_{abcd}\,e^a\wedge e^b\wedge e^c\wedge e^d
+
\int e\,\mathcal{L}_{\mathrm{SM}}.
\label{eq:nsm}
\end{equation}
In the certified regime the microscopic description is a
holographic pair \cite{Maldacena1998,Witten1998}: a unitary CFT on
the conformal boundary of AdS, with a bulk dual whose infrared is
\eqref{eq:nsm}. Two readings must not be mixed.

\begin{definition}[Reading A, the monograph]
The Standard Model fields in \eqref{eq:nsm} are \emph{bulk} fields.
The boundary CFT is some other theory, with a global flavour
symmetry whose bulk gauge fields are $G_{\mathrm{SM}}$.
\end{definition}

\begin{definition}[Reading B, the temptation]
The Standard Model \emph{is} the boundary CFT, and bulk Einstein--Cartan
is its entanglement geometry.
\end{definition}

Reading~B would be a selection. It is false, by the next section.
Reading~A is the monograph's reading. Under it, Q4a is: which CFT?

\section{The Standard Model is not a holographic CFT}

A necessary condition for a four-dimensional gauge theory to be a
CFT is the vanishing of all one-loop beta-function coefficients in
the convention $\beta_i=b_i g_i^3/(16\pi^2)$. From the SM
representation table, with $Q=T_3+Y/2$ and GUT-normalized $U(1)$,
\begin{equation}
b_1=\frac{41}{10},
\qquad
b_2=-\frac{19}{6},
\qquad
b_3=-7.
\label{eq:b}
\end{equation}
All three are nonzero. The Standard Model is not a CFT. \pP{}
(\emph{computed} from Dynkin indices; textbook values recovered).

A necessary condition for a \emph{classical} Einstein dual is a
parametrically large $N$ and a sparse spectrum of light operators
with a large gap \cite{Heemskerk2009}. Before electroweak symmetry
breaking the SM has $118$ on-shell light degrees of freedom
($24$ gauge helicities, $4$ Higgs, $90$ fermionic). That is
$O(10^2)$, not a limit $N\to\infty$. \pP{} (count).

\begin{theorem}[Reading B is false]
\label{thm:B}
The Standard Model is not the boundary theory of a holographic pair
with a classical Einstein--Cartan bulk. \pP{} (necessary conditions
fail).
\end{theorem}

A mutation confirms the CFT check is not tautological: $\mathcal{N}=4$
super-Yang--Mills has one-loop $b=0$ by the adjoint index identity
$T(\mathrm{adj})=C_2(G)$. The same script that rejects the SM accepts
that necessary condition for $\mathcal{N}=4$. Accepting the necessary
condition is not a selection of $\mathcal{N}=4$ as the pair.

\section{What the certified principles actually see}

The Casini--Huerta--Myers modular Hamiltonian of a ball of radius
$R$ in a CFT vacuum is \cite{CHM2011}
\begin{equation}
H_B
=
2\pi\int_{\lvert x\rvert<R}\!d^{d-1}x\,
\frac{R^2-\lvert x\rvert^2}{2R}\,T_{00}(x).
\label{eq:chm}
\end{equation}
Faulkner--Guica--Hartman--Myers--Van Raamsdonk convert the first law
$\delta S_B=\delta\langle H_B\rangle$ into linearized Einstein
equations sourced by $\delta\langle T_{\mu\nu}\rangle$
\cite{FGHMV2014}. The only CFT data in that argument are the stress
tensor, the central charge that normalizes it, and the ball.

Paper~II selects the Young symmetry of an explicit spin-current
source: mixed-symmetry non-minimal couplings are relevant
deformations that destroy the conformal frame; the surviving sector
is axial and minimal \cite{McGwierII}. That is a statement about
\emph{how} a spinor couples to $\omega$, not about
\emph{which} gauge group the spinor is charged under.

\begin{theorem}[Underdetermination]
\label{thm:under}
Let $\mathcal{P}$ be the set of holographic pairs whose bulk
infrared contains Einstein gravity, a cosmological term, and some
minimally coupled anomaly-free chiral gauge theory with an axial
current. The maps
\begin{align*}
G_{\mathrm{SM}}
&\;\longrightarrow\;
G',
\qquad
n_g
\;\longrightarrow\;
n_g',
\qquad
Y_{u,d,e}
\;\longrightarrow\;
Y',
\end{align*}
stay inside $\mathcal{P}$ whenever $G'$ is compact and the new
spectrum is anomaly-free. None of $G_{\mathrm{SM}}$, $n_g$, or the
Yukawas appears in \eqref{eq:chm} or in the FGHMV source. The
certified principles therefore do not select a unique element of
$\mathcal{P}$. \pP{} (invariance).
\end{theorem}

Theorem~\ref{thm:under} is the parameter-count test of this
question. The first law is one structural constraint on the bulk
metric. The matter ledger of the NSM is a finite list of installed
choices. The space of holographic CFTs is infinite-dimensional
(spectra, OPE coefficients, flavour groups). $N_{\mathrm{free}}$
exceeds $N_{\mathrm{constraint}}$. Uniqueness fails.

\section{What \emph{is} selected}

Given a holographic pair, the certified principles do select the
bulk gravitational sector at linear order about AdS: the metric
equation is Einstein \cite{FGHMV2014}, and the torsion coupling,
if spinors are present, is the axial Einstein--Cartan one of
Paper~II (with the axial obstruction of Paper~14 standing at
second order). That is ``certify, select, assemble'' applied to
gravity, which is what Papers~I--III already claimed.

Q4a asked for the same operation on the pair itself. The operation
does not extend. Gravity is selected \emph{given} a pair. The pair
is not selected.

\section{Existence remains open}

Bottom-up, \eqref{eq:nsm} in AdS \emph{is} a holographic model. That
is the same effective field theory, not a microscopic pair.

Top-down, a string or M-theory vacuum whose infrared is Palatini
Einstein--Cartan plus the Standard Model is not constructed here and
is not selected by entanglement first-principles. Citing an existing
compactification as ``the'' pair would be a fabrication of
uniqueness.

Paper~15's quadratic axial deformation is a change of the bulk
ultraviolet, not a choice of CFT.

\begin{admission}
I did not select a holographic pair. I proved that the program's
own principles cannot select one, and that the Standard Model
cannot be one. Existence of some pair whose bulk infrared is the
NSM is still \pO{}. That is the proper-context answer to Q4a.
\end{admission}

\section{Machine check}

Solver \texttt{m9\_5\_q4a\_pair.py}; auditor
\texttt{m9\_5\_audit\_q4a.py} (species table, no solver import).
Gates C1--C5 PASS. The auditor recovers $(b_1,b_2,b_3)$ by summing
representations generation-by-generation. C3 is structural (the
CHM integrand), not a numerical search.

\section{Verdict}

\begin{center}
\begin{tabular}{@{}p{0.48\textwidth}p{0.44\textwidth}@{}}
\toprule
Claim & Status \\
\midrule
Q4a as ``invent a CFT'' & refused (Paper~11) \\
Reading B (SM is the boundary CFT) & false, Thm.~\ref{thm:B} \\
Selection-uniqueness of the pair & false, Thm.~\ref{thm:under} \\
SM one-loop $(b_1,b_2,b_3)$ & computed, audited \\
Bulk gravity given a pair & already certified, Papers~I--III \\
Existence of some pair & still \pO{} \\
Final Status ``UV completion'' & unchanged, \pO{} \\
\bottomrule
\end{tabular}
\end{center}

The phrase that is true: \emph{Q4a, asked as a first-principles
selection of the holographic pair, has a negative answer.}
The phrase that is false: \emph{we have the ultraviolet completion.}

\begin{thebibliography}{99}
\bibitem{Maldacena1998}
J.~Maldacena,
Adv.\ Theor.\ Math.\ Phys.\ 2 (1998) 231, arXiv:hep-th/9711200.
\bibitem{Witten1998}
E.~Witten,
Adv.\ Theor.\ Math.\ Phys.\ 2 (1998) 253, arXiv:hep-th/9802150.
\bibitem{CHM2011}
H.~Casini, M.~Huerta and R.C.~Myers,
JHEP 05 (2011) 036, arXiv:1102.0440.
\bibitem{FGHMV2014}
T.~Faulkner, M.~Guica, T.~Hartman, R.C.~Myers and M.~Van Raamsdonk,
JHEP 03 (2014) 051, arXiv:1312.7856.
\bibitem{Heemskerk2009}
I.~Heemskerk, J.~Penedones, J.~Polchinski and J.~Sully,
JHEP 10 (2009) 079, arXiv:0907.0151.
\bibitem{McGwierII}
R.W.~McGwier,
``Spin-Current Sources and Ball Modular Hamiltonians,''
\url{https://github.com/n4hy/New_Model_Emergent_Gravity}.
\bibitem{McGwierMonograph}
R.W.~McGwier,
``Emergent Gravity and a New Lagrangian,'' ibid.
\bibitem{McGwierXI}
R.W.~McGwier,
``Open Questions, Attempts, and Admissions,'' ibid.
\end{thebibliography}

\end{document}
